\section*{Hoofdstuk \ref{ch:b2:diatomic-mos} --- Molecuulorbitalen van tweeatomige moleculen}

\begin{solution}{exo:b2:diatomic-mos:1}
$E_+ = (-13.6 - 9.0)/1.60 = \qty{-14.1}{eV}$; $E_- = (-13.6 + 9.0)/0.40 = \qty{-11.5}{eV}$.
Twee elektronen in $\psi_+$: $-28.25$ tegenover $2 \times (-13.6) = \qty{-27.2}{eV}$, in dit
ruwe model gebonden met \qty{1.05}{eV}.
\end{solution}

\begin{solution}{exo:b2:diatomic-mos:2}
\ce{He2}: $(2 - 2)/2 = 0$; \ce{He2+}: $(2 - 1)/2 = 0.5$; \ce{Li2}: 1; \ce{B2}: zes valentie-elektronen,
$\sigma_{2s}^2\sigma_{2s}^{*2}\pi_{2p}^2$, \omterm{def:b2:diatomic-mos:bond-order}{bindingsorde} 1.
\end{solution}

\begin{solution}{exo:b2:diatomic-mos:3}
\ce{B2} (twee ongepaarde elektronen in de twee $\pi_{2p}$-orbitalen, met s-p-menging) en
\ce{O2} (twee in $\pi^*$). \ce{C2}, \ce{N2}, \ce{F2} zijn \omterm{def:b2:diatomic-mos:magnetism}{diamagnetisch}.
\end{solution}

\begin{solution}{exo:b2:diatomic-mos:4}
Verschillend van nul: (1s, $2\mathrm p_z$), ($2\mathrm p_x$, $2\mathrm p_x$), (2s, $2\mathrm p_z$).
Nul door symmetrie: (1s, $2\mathrm p_x$) en ($2\mathrm p_x$, $2\mathrm p_y$).
\end{solution}

\begin{solution}{exo:b2:diatomic-mos:5}
Ionisatie verwijdert een bindend elektron: \omterm{def:b2:diatomic-mos:bond-order}{bindingsorde} 2.5 in plaats van 3, dus een langere en
zwakkere binding, zoals $D_0(\ce{N2+}) = \qty{8.71}{eV} < D_0(\ce{N2}) = \qty{9.76}{eV}$ laat zien.
\end{solution}

\begin{solution}{exo:b2:diatomic-mos:6}
De orbitalen van zuurstof liggen lager (zuurstof is elektronegatiever): volgens
\cref{thm:b2:diatomic-mos:heteronuclear} liggen de \omterm{def:b2:diatomic-mos:bonding}{bindende orbitalen} zwaarder op het
lagere atoom, zuurstof, en de hoger gelegen orbitalen op koolstof. De hoogste bezette
orbitaal, een $\sigma$-orbitaal met overwegend koolstofkarakter die van zuurstof weg wijst, is
het vrije elektronenpaar op koolstof.
\end{solution}

\begin{solution}{exo:b2:diatomic-mos:7}
Alleen de $2\mathrm p_z$ van fluor (langs de as) heeft de symmetrie van de 1s van waterstof:
samen vormen ze een $\sigma$-orbitaal die zwaarder op F ligt en een $\sigma^*$ die zwaarder op H ligt. De
$2\mathrm p_x$, $2\mathrm p_y$ van fluor (overlap nul) en zijn lage 2s (te ver in energie)
blijven niet-bindend: drie vrije elektronenparen. De acht valentie-elektronen vullen 2s, $\sigma$
en de twee niet-bindende p.
\end{solution}

\begin{solution}{exo:b2:diatomic-mos:8}
$\bar\alpha = \qty{-12}{eV}$, $\Delta = \qty{4}{eV}$, $\sqrt{4 + 9} = \qty{3.61}{eV}$:
$E = \qty{-15.6}{eV}$ en \qty{-8.4}{eV}. Voor het lagere niveau is $(\alpha_B - E)c_B = -\beta
c_A$: $1.61c_B = 3.0c_A$, $c_A/c_B = 0.535$, $c_B^2 = 1/(1 + 0.286) = 0.78$.
\end{solution}

\begin{solution}{exo:b2:diatomic-mos:9}
Acht valentie-elektronen: $\sigma_{2s}^2\sigma_{2s}^{*2}\pi_{2p}^4$, waarbij de orbitaal $\sigma_{2p}$
(door de menging boven $\pi_{2p}$ geduwd) leeg is. \omterm{def:b2:diatomic-mos:bond-order}{Bindingsorde} $(6 - 2)/2 = 2$, afkomstig van de
twee gevulde $\pi$-orbitalen.
\end{solution}

\begin{solution}{exo:b2:diatomic-mos:10}
NO, elf valentie-elektronen, waarvan één in $\pi^*$: \omterm{def:b2:diatomic-mos:bond-order}{bindingsorde} 2.5. \ce{NO+}, tien: \omterm{def:b2:diatomic-mos:bond-order}{bindingsorde}
3. $D_0(\ce{NO+}) = 6.51 + 13.62 - 9.26 = \qty{10.87}{eV}$: veel sterker dan NO, omdat
het verwijderde elektron antibindend was.
\end{solution}

\begin{solution}{exo:b2:diatomic-mos:11}
$E_+ + E_- = \dfrac{(\alpha + \beta)(1 - S) + (\alpha - \beta)(1 + S)}{1 - S^2} = \dfrac{2(\alpha -
\beta S)}{1 - S^2}$. Dit is groter dan $2\alpha$ wanneer $\alpha - \beta S > \alpha - \alpha S^2$, dus
$\beta < \alpha S$ (delen door $-S < 0$). Met vier elektronen, twee in elke orbitaal, ligt de
totale energie boven die van de gescheiden orbitalen: twee gevulde orbitalen stoten elkaar af.
\end{solution}

\begin{solution}{exo:b2:diatomic-mos:12}
\omterm{def:b2:diatomic-mos:bond-order}{Bindingsordes} 2.5, 2, 1.5, 1: in die volgorde worden de bindingen langer en zwakker. Waterstofperoxide
bevat de enkelvoudige O--O-binding van peroxide (eenheid \ce{O2^2-}); \ce{KO2} bevat
het superoxide-ion \ce{O2-}, \omterm{def:b2:diatomic-mos:magnetism}{paramagnetisch}.
\end{solution}

\begin{solution}{pb:b2:diatomic-mos:1}
\textbf{1.} Twaalf.
\textbf{2.} $\sigma_{2s}^2\,\sigma_{2s}^{*2}\,\sigma_{2p}^2\,\pi_{2p}^4\,\pi_{2p}^{*2}$, met de twee $\pi^*$-elektronen
in verschillende orbitalen.
\textbf{3.} $(8 - 4)/2 = 2$.
\textbf{4.} Twee: \ce{O2} is \omterm{def:b2:diatomic-mos:magnetism}{paramagnetisch}.
\textbf{5.} Ze paart alle elektronen, terwijl er twee alleen zitten in twee ontaarde
$\pi^*$-orbitalen.
\textbf{6.} $D_0 = 2 \times 246.79 = \qty{493.6}{kJ/mol}$, \qty{5.12}{eV}.
\textbf{7.} \ce{O2+}: $\pi^{*1}$, \omterm{def:b2:diatomic-mos:bond-order}{bindingsorde} 2.5.
\textbf{8.} \ce{O2-}: $\pi^{*3}$, \omterm{def:b2:diatomic-mos:bond-order}{bindingsorde} 1.5.
\textbf{9.} \ce{O2^2-}: $\pi^{*4}$, \omterm{def:b2:diatomic-mos:bond-order}{bindingsorde} 1.
\textbf{10.} \ce{O2+} één, \ce{O2-} één, \ce{O2^2-} geen.
\textbf{11.} \ce{O2+} $<$ \ce{O2} $<$ \ce{O2-} $<$ \ce{O2^2-}.
\textbf{12.} Het peroxide-ion.
\textbf{13.} \ce{O2+}, \ce{O2} en \ce{O2-}.
\textbf{14.} Uit een $\pi^*$-orbitaal.
\textbf{15.} De $\pi^*$-orbitaal, antibindend, ligt boven het 2p-niveau van het atoom: haar
elektron is minder sterk gebonden.
\textbf{16.} Twee wegen van \ce{O2} naar $\ce{O} + \ce{O+} + \mathrm e^-$: dissociëren en daarna
een atoom ioniseren, $D_0(\ce{O2}) + I(\ce{O})$; of het molecuul ioniseren en daarna het ion dissociëren,
$I(\ce{O2}) + D_0(\ce{O2+})$. Gelijkstellen: $D_0(\ce{O2+}) = 5.12 + 13.62 - 12.07 = \qty{6.66}{eV}$.
\textbf{17.} Groter dan $D_0(\ce{O2}) = \qty{5.12}{eV}$: een antibindend elektron verwijderen
versterkt de binding (orde 2.5).
\textbf{18.} $D_0(\ce{N2}) = 2 \times 470.82/96.485 = \qty{9.76}{eV}$; $D_0(\ce{N2+}) = 9.76 +
14.53 - 15.58 = \qty{8.71}{eV}$.
\textbf{19.} De hoogste bezette orbitaal van \ce{N2} is bindend, onder het 2p-niveau van het
atoom.
\textbf{20.} Wanneer het verwijderde elektron uit een \omterm{def:b2:diatomic-mos:bonding}{antibindende orbitaal} komt.
\textbf{21.} De regel van Hund (maximaal aantal parallelle spins in ontaarde orbitalen).
\textbf{22.} Beide $\pi^*$-elektronen in dezelfde orbitaal, met gepaarde spins.
\textbf{23.} Twee elektronen in dezelfde orbitaal stoten elkaar sterker af, en verliezen de
uitwisselingsstabilisatie van parallelle spins.
\textbf{24.} \omterm{def:b2:diatomic-mos:bond-order}{Bindingsorde} nog steeds 2; alle elektronen gepaard: \omterm{def:b2:diatomic-mos:magnetism}{diamagnetisch}.
\textbf{25.} \ce{O2+} heeft \omterm{def:b2:diatomic-mos:bond-order}{bindingsorde} $\boldsymbol{2.5}$ en een dissociatie-energie van
$\boldsymbol{\approx \qty{6.66}{eV}}$, tegenover 2 en \qty{5.12}{eV} voor \ce{O2}.
\end{solution}
